% ============================================================================
% Appendix E.1: Limitations and Applicability
% ============================================================================
% Purpose: Clarify scope, generalizability, and when findings apply

\subsection{Limitations and Applicability}
\label{sec:limitations}

To establish the scope and generalizability of our approach, we systematically
evaluated our architecture under varying conditions. Several important
considerations emerged from this analysis.

% ----------------------------------------------------------------------------
% Limitation 1: Domain Mismatch Severity
% ----------------------------------------------------------------------------
\paragraph{Prior Quality Dependency.}
Our evaluation focuses on a challenging deployment scenario: severe domain
mismatch where 68.6\% of warmup training prompts are classified as ``hard,''
but only 13.7\% of holdout prompts fall into this category.
This represents a substantial distribution shift where warmup priors are
systematically misspecified.
Under these conditions (Figure~3), constant exploration ($\alpha=2.0$) prevents
premature commitment to incorrect beliefs, achieving 48\% improvement over
adaptive decay (60.6 vs 90.2 regret).

However, this finding may not generalize to scenarios with \emph{accurate}
priors and \emph{low} domain mismatch.
When warmup data closely matches the deployment distribution, adaptive
exploration schedules (e.g., $\alpha_t = \alpha_0/\sqrt{t}$) may improve
sample efficiency by reducing unnecessary exploration once the system has
learned accurate model preferences.
Our results suggest that the optimal exploration schedule depends critically
on prior quality---a dimension not traditionally considered in bandit algorithm
design.

% ----------------------------------------------------------------------------
% Limitation 2: Strategy Overhead
% ----------------------------------------------------------------------------
\paragraph{Strategy Selection Trade-offs.}
While Corralling provides robust performance across diverse scenarios
(59.2 regret), our experiments reveal that practitioners with \emph{validated}
prior quality can achieve superior performance using simpler single-expert
strategies.
Specifically, pure Tabula Rasa achieves 49.5 regret when priors are known to
be severely misspecified---16\% better than Corralling's hedging overhead.

This finding challenges the assumption that meta-learning approaches are
universally optimal.
The 20\% overhead of hedging is justified only under \emph{genuine} uncertainty
about prior quality.
In production deployments where prior quality can be validated on held-out
deployment data (N=100--200 samples), practitioners should consider strategy
selection based on validation results rather than defaulting to Corralling.

The value of Corralling lies in its \emph{safety guarantee}: 18.5\%
improvement over harmful warmup-only deployment (59.2 vs 74.7 regret).
This safety property is most valuable in risk-averse deployments where prior
quality is truly uncertain and the cost of poor performance is high.

% ----------------------------------------------------------------------------
% Limitation 3: Heterogeneous Alpha Hypothesis
% ----------------------------------------------------------------------------
\paragraph{Mechanism Validation Through Ablation.}
Our ablation study (Appendix~\ref{appendix:mathematical_foundations},
Table~\ref{tab:corralling_ablation}) validated that the system's robustness
in the Figure~3 mismatch scenario stems from \emph{constant} exploration
($\alpha=2.0$ for both experts) rather than \emph{heterogeneous} exploration
schedules.
Empirical validation revealed that homogeneous constant exploration
outperforms mixed schedules by 6.3\%.
This demonstrates that the system's effectiveness derives from hedging between
\emph{different initialization sources} (semantic priors vs.\ online learning)
while maintaining consistent vigilance, rather than from diversified
exploration strategies.

The ablation study establishes that the simpler homogeneous design provides
superior performance while reducing system complexity.  Note, however, that the
ablation grid covers 15 configurations with 3 seeds each
(Remark~\ref{remark:ablation_scope}); broader grids may reveal additional
interactions.

% ----------------------------------------------------------------------------
% Limitation 4: Variance and Seed Dependency
% ----------------------------------------------------------------------------
\paragraph{Outcome Variance and Reproducibility.}
Our multi-seed validation reveals substantial variance in expert weight
evolution: final warmup weight is 0.382 $\pm$ 0.471 (standard deviation
exceeds mean).
This high variance indicates that outcomes are seed-dependent and influenced
by the stochastic ordering of requests.
Practitioners cannot rely on deterministic predictions (e.g., ``weight will be
0.2 after 200 requests'') and must instead monitor their specific deployment's
weight trajectory.

This variance also has implications for experimental methodology: single-seed
evaluations are insufficient for characterizing bandit algorithm performance.
Our main experiments use 20 seeds per configuration; the ablation grid
(Table~\ref{tab:corralling_ablation}) uses 3 seeds per cell.
While the ablation provides directional evidence, variance-sensitive
conclusions (e.g., precise configuration rankings) would benefit from larger
sample sizes.

% ----------------------------------------------------------------------------
% Limitation 5: Regime-Dependent Hyperparameter Effects
% ----------------------------------------------------------------------------
\paragraph{Regime-Dependent Behavior.}
Multi-seed validation reveals two distinct operating regimes: a
warmup-dominant regime (approximately 33\% of data orderings) where Corralling
maintains the semantic transfer expert, and a tabula rasa-dominant regime
(approximately 67\%) where it abandons semantic transfer in favor of cold-start
exploration.
This regime dependence is a feature of Corralling's adaptive capacity: when
warmup priors match the deployment distribution, hyperparameter choice
($\neff$) significantly affects performance; when priors mismatch, Corralling
correctly detects over-confidence and switches strategies, rendering
$\neff$ choice less consequential.
The specific regime frequencies observed here reflect our LMSYS Arena test
distribution and may differ for other model portfolios or task distributions.

% ----------------------------------------------------------------------------
% Limitation 6: Computational Considerations
% ----------------------------------------------------------------------------
\paragraph{Computational Overhead.}
While our coordinator adds negligible latency ($<$1ms) and modest memory
overhead (2$\times$ for maintaining dual experts), practitioners deploying at
scale should consider:
(1)~Expert weight logging incurs storage costs;
(2)~Multi-expert inference requires maintaining separate confidence bounds;
(3)~Importance-weighted loss computation requires tracking selection
probabilities.
For extremely latency-sensitive applications, these overheads may be
non-negligible.

% ----------------------------------------------------------------------------
% Generalizability
% ----------------------------------------------------------------------------
\paragraph{Generalizability.}
Our evaluation uses a single dataset (LMSYS Arena) with two models
(Mixtral-8x7B, GPT-4-Turbo) under a specific mismatch pattern (chatbot arena
$\to$ holdout distribution).
While this scenario is representative of real-world LLM deployment challenges,
additional validation on diverse datasets, model pairs, and mismatch patterns
would strengthen generalizability claims.
The rapid adaptation observed in our experiments (16 requests) may vary with
mismatch severity: mild mismatch might require more samples for detection,
while catastrophic mismatch might be even faster.

% ----------------------------------------------------------------------------
% Summary
% ----------------------------------------------------------------------------
Despite these limitations, our core contributions remain robust:
(1)~Constant exploration prevents premature exploitation under mismatch;
(2)~Corralling provides safety when prior quality is uncertain;
(3)~Strategy selection based on validated prior quality optimizes performance;
(4)~Rigorous ablation studies are essential for validating architectural
hypotheses.
These insights provide actionable guidance for practitioners deploying
adaptive LLM routing systems in uncertain environments.
